\documentclass[11pt,a4paper]{article}
\usepackage[margin=27mm]{geometry}
\usepackage{amsmath,amssymb,amsthm,hyperref}
\emergencystretch=1em
\newtheorem{lemma}{Lemma}\newtheorem{proposition}[lemma]{Proposition}
\newcommand{\Sz}{\operatorname{Sz}}
\title{Diagrams of the surviving prescribed families}
\author{Human Compression project}\date{}
\begin{document}\maketitle
\begin{abstract}
The diagram remarks in the fixed higher-rank Suzuki source concern
flag-transitive realizations of its prescribed intersection lattices.
A subgroup-product test verifies those remarks for the surviving
families. The test does not assume the source's exhaustive list and
makes no additional assertion of flag transitivity.
\end{abstract}

\section{A product test}
\begin{lemma}\label{product}
Let a flag-transitive coset geometry have standard rank-two residue
$\Gamma(H;A,B)$. This residue is a generalized digon if and only if
$H=AB$, equivalently $|H|\,|A\cap B|=|A|\,|B|$.
\end{lemma}
\begin{proof}
All objects of the two types are incident precisely when every left
$A$-coset meets every left $B$-coset. By translation this is equivalent
to $A$ meeting $hB$ for every $h\in H$, hence to $h\in AB$ for
every $h$. The product has cardinality $|A||B|/|A\cap B|$.
The geometry diagram omits exactly the pairs with a digon residue.
\end{proof}

Write $e\ge3$ odd, $t=2^e$, and let $a_1,\ldots,a_k$ be distinct
odd primes. Put $n=e\prod_i a_i$ and $S_d=\Sz(2^d)$.
In all families below the field parabolics are $P_i=S_{n/a_i}$ in
the acquired simultaneous position. The intersection formulas proved
in the companion manuscripts apply to every subset of the types.

\section{The families with a torus parabolic}
\begin{proposition}
For the canonical realizations of Constructions 8 and 9, the diagram
is a star with centre the torus-normalizer type $K$ and leaves the
$k$ field types.
\end{proposition}
\begin{proof}
Use the notation $\gamma$, $\ell_i$ and $c\in\{2,4\}$ of the torus
family theorem. Thus $K=C_{\gamma\prod_i\ell_i}:C_c$, and
$K\cap S_J=C_{\gamma\prod_{i\notin J}\ell_i}:C_c$.
Flag transitivity has already been proved by descent and cyclic coset
coordinates, so the standard subgroup residues are the actual residues.

For two field types $i,j$, all other types, including $K$, have been
selected. The current group has order $c\gamma\ell_i\ell_j$;
its two parabolics have orders $c\gamma\ell_i,c\gamma\ell_j$,
and their intersection has order $c\gamma$. Lemma~\ref{product}
makes this residue a digon.

For the pair $i,K$, the current group is $\Sz(t^{a_i})$;
the parabolics are $\Sz(t)$ and $C_{\gamma\ell_i}:C_c$,
with intersection of order $c\gamma$. Their product has size
$|\Sz(t)|\ell_i$. Its two-part is $t^2$, whereas the two-part
of $|\Sz(t^{a_i})|$ is $t^{2a_i}$. They cannot be equal.
Every field type is therefore adjacent to $K$, and no two field
types are adjacent to one another.
\end{proof}

\section{The families with two soluble parabolics}
The corrected Constructions 5 and 7 use
\[
 M=(t-1)\prod_i\ell_i,\quad
 P_D=C_M:C_2,\quad P_E=Z_{2^n}:C_M,
\]
where the split-torus scalar subgroup has order $M$. These subgroups
and every intersection have been constructed in the companion scalar
field theorem. That theorem proves the source's subgroup maximality
condition, but does not prove flag transitivity.
\begin{proposition}
If a realization with these prescribed intersections is a
flag-transitive geometry, its diagram is a star with centre $D$;
its leaves are $E$ and all the field types.
\end{proposition}
\begin{proof}
Put $\gamma=t-1$. For two field types $i,j$, after selecting both
$D,E$ and all other fields, the current group is cyclic of order
$\gamma\ell_i\ell_j$. Its two parabolics have orders
$\gamma\ell_i,\gamma\ell_j$, with meet of order $\gamma$.
Their product is the entire group.

For types $i,E$, the current group has order $2\gamma\ell_i$;
the two parabolics have orders $2\gamma$ and $\gamma\ell_i$,
and their meet has order $\gamma$. Again the residue is a digon.
For types $i,D$, the current group has order
$t^{a_i}\gamma\ell_i$; its parabolics have orders
$t\gamma$ and $\gamma\ell_i$, with meet $\gamma$.
Their product has order $t\gamma\ell_i$, strictly smaller.
Finally, for $D,E$ the current group is $\Sz(t)$, and the product
of its parabolics has order $2t\gamma$, smaller than
$t^2(t^2+1)\gamma$. These are exactly the asserted edges.
\end{proof}

\section{Pure subfield geometries}
\begin{proposition}
The canonical pure-subfield geometry has complete diagram on its
$k$ field types. In particular its diagram is not a star when $k>3$.
\end{proposition}
\begin{proof}
For a pair of remaining relative primes $a,b$, the current group
is $\Sz(t^{ab})$, its parabolics are $\Sz(t^a),\Sz(t^b)$,
and their intersection is $\Sz(t)$. Here the bottom $t=2^e$
may equal two; the group-order formula remains valid then.
Write $f(x)=x^2(x^2+1)(x-1)$. For $x\ge2$,
\[
 \tfrac12 x^5<f(x)<x^5.
\]
Consequently
\[
 \frac{f(t^{ab})f(t)}{f(t^a)f(t^b)}
 >\frac{t^{5(ab+1-a-b)}}4>1,
\]
since $a,b\ge3$ and $t\ge2$. The product criterion fails for
every pair, as claimed. Flag transitivity and the residue identification
were proved in the pure-subfield theorem.
\end{proof}

\paragraph{Source scope and attribution.}
D.~Leemans, \emph{The residually weakly primitive pre-geometries of the
Suzuki simple groups}, Note di Matematica 20(1) (2000/2001), pp.~18--19,
observes that the listed higher-rank families other than Construction~11
have star diagrams, conditional on being flag-transitive geometries.
The preceding proofs establish that statement for the surviving
prescribed families. The claimed pentagon-and-circles diagram of
Construction~1 is conditional on an intersection lattice now proved
impossible, and therefore supplies no existing geometry. Constructions
2--4 and 6 have the same obstruction to their prescribed lattices;
Construction~10 fails the source maximality axiom. None of these
conditional diagram calculations restores the false exhaustive list.
The additional statement that almost all Suzuki geometries have this
form has no specified counting measure and is not a classification
conclusion established here. The source's future request to study flag
transitivity is answered for the canonical families 8, 9 and 11;
flag transitivity of arbitrary embeddings and of the corrected families
5 and 7 is not asserted by this manuscript.
\end{document}
